\documentclass[12pt,a4paper]{article}
\usepackage[scheme=plain]{ctex}
\usepackage[margin=2.5cm]{geometry}
\usepackage{indentfirst}
\usepackage{setspace}
\usepackage{fontspec}
\setCJKmainfont[BoldFont=SimHei,ItalicFont=KaiTi]{SimSun}
\newCJKfontfamily\kaifont{KaiTi}[AutoFakeBold]
\usepackage{newpxtext,newpxmath}

\usepackage{lastpage}





\pagestyle{plain}

\usepackage[nobottomtitles*]{titlesec}
\titleformat{\section}{\Large\bfseries}{}{0em}{ }
\titlespacing*{\section}{0pt}{3ex}{1.5ex}

\usepackage{xcolor}
\usepackage{needspace}

\setlength{\parindent}{2em}
\setlength{\parskip}{0.4em}

\doublespacing

\title{\bfseries Day 4\\[0.3em]共谋、非正式制度与组织行为}
\author{}
\date{}

\begin{document}
\maketitle

\vspace{-2em}

\vspace{1em}


\vspace{-0.3em}{\color{gray}\small\kaifont \noindent 理解组织行为有两种常见路径。一是看正式制度，谁管谁，什么事走什么流程。二是看个人素质，这个人是不是偷懒了，那个人是不是搞腐败了。两种路径都有解释力，但都漏掉了一类现象。有些行为既不是规则上写着的，也不归因于哪个人的错误。它稳定存在，反复再生，组织中的所有人都知道它在发生，也都默认它是正常的。周雪光把这类现象称为“制度化了的非正式行为”。他以中国基层政府间“上有政策、下有对策”的共谋现象为切入口，论证了一个在组织研究中具有普遍意义的命题。正式制度本身会催生出大量非正式行为，这些行为不是制度的例外，而是制度的实现形式。}\vspace{0.5em}
\vspace{0.3em}

\noindent\rule{\textwidth}{0.4pt}

\vspace{0.3em}



[内容提要]中国政府行为的一个突出现象是，在执行来自上级部门特别是中央政府的各种指令政策时，基层上下级政府常常共谋策划、暗渡陈仓，采取“上有政策、下有对策”的各种手段，来应付这些政策要求以及随之而来的各种检查，导致了实际执行过程偏离政策初衷的结果。本文从组织学角度，对这类现象提出一个理论解释。本文的中心命题是：在中国行政体制中，基层政府间的共谋行为已经成为一个制度化了的非正式行为；这种共谋行为是其所处制度环境的产物，有着广泛深厚的合法性基础。本文讨论政府组织制度的三个悖论，对这一现象提出理论解释：一是政策一统性与执行灵活性的悖论，二是激励强度与目标替代的悖论，三是科层制度非人格化与行政关系人缘化的悖论。本研究强调，共谋行为不能简单地归咎于政府官员或执行人员的素质或能力，其稳定存在和重复再生是政府组织结构和制度环境的产物；是现行组织制度中决策过程与执行过程分离所导致的结果；在很大程度上也是近年来政府制度设计特别是集权决策过程和激励机制强化所导致的非预期结果。而欲改变这一状况，首先需要对政府组织现象进行深入系统的研究，提出有力的理论解释。



一、现象与问题



作为中央政策执行过程的最终环节，基层政府在中国政府组织制度中有着极为重要的战略位置。已有许多研究工作关注分析基层政府的状况和行为（参见赵树凯，2005，2006b；在社会学领域，参见周飞舟，2006；张静，2000；吴毅，2007）。在中国政府组织制度中，贯彻执行上级政府的各项政策指令是基层政府的一个重要工作内容。上级决策者和政策研究者经常用“上传下达、令行禁止”这样的词汇来形容基层政府工作的理想状态。但实际生活中一个引人注目的现象是，在执行来自上级部门特别是中央政府的各种政策指令时，一些基层政府常常共谋策划，采取“上有政策、下有对策”的各种做法，联手应付自上而下的政策要求以及随之而来的各种检查，导致了实际执行过程偏离政策初衷的结果。例如，一位镇政府主管计生工作的官员这样描述有关的检查工作：



\begin{quotation}\kaifont\noindent 省市县都组织检查。当省里来检查时，市、县、镇联合起来对付省的检查团；当市里来检查时，县、镇就联合起来对付市里的检查团。省里到县里检查时，事先不通知去哪个镇或村。但是，实际上各地政府都联合起来对付检查团。省里来的时候，市里、县里就提前来打招呼：“你们那里的问题处理好……”检查团到了县里，各个镇都得到通知，严阵以待。一旦领队的得到通知去哪个村，在去那个村的路上就有电话通知到那个村，连车号、时间、地点的消息都十分详尽。检查团一般都是早晨8点以前就到达。所以，一大早，村里妇女小组长就在各个路口把守，一旦发现情况就立即通知村里，把有问题的孩子转移出去……我工作6年来，除了有一个乡被省检查团查出来外，只有两次（两个镇）是县里检查出计划外的问题，做了一票否决的处理，镇书记和镇长都调换工作，没有重用。县里也没有向上报告。（访谈记录 H2346）\end{quotation}



这种上下级政府间联手应对更上级政府的现象普遍存在。媒介报道更多地注意到了这些负面状况。例如，煤窑安全事故发生后，事发地基层政府间层层包庇，封锁道路，某些基层官员滥用权力，压迫民众；一旦事发，自上而下调查时，那些基层政府上下级也经常互相掩盖，大事化小，小事化了（参见陈桂棣、春桃，2004；高王凌，2006）。这些行为有着各种表现形式，或者是基层政府出现问题而游说直接上级部门为之掩盖保护；或者是直接上级部门要求下级部门掩盖问题，以应付来自更上级政府的要求和检查。更多的情形下，这类行为体现在上下级基层政府持续双向的互动过程中。因为这些应对策略和行为常常与上级政策指令不符甚至相悖，所以它们大多是通过非正式方式加以实施（例如，“打招呼”、私下安排等）。本文用基层政府间“共谋行为”这一概念来概括描述这类非正式行为。



值得强调的是，本文所关注的基层政府“共谋行为”是在中性意义上使用的一个分析概念。如下文所讨论的，这些“共谋现象”在很多情况下是基层政府执行政策时所表现的灵活性和适应性，是制度创新的一个重要源泉。而那些负面意义上的基层政府间共谋行为，虽然是中央政府所不愿意看到并且力图克服的，但这类行为的范围之广、程度之深、表现之公开、运作之坚韧却是少见的。在中国社会转型过程中，随着政府制度职能的演变和改革，这一问题也随之日益突出，甚至有愈演愈烈的趋势。重复再现的组织现象是建筑在稳定持续的组织制度基础之上和相应的组织环境之中的。本文对这一现象从组织学的角度尝试提出一个理论解释。



我们的中心命题是：在中国行政体制中，基层政府间的共谋行为已经成为一个制度化了的非正式行为。即，这种共谋行为是基层政府所处制度环境的产物，因此有着广泛深厚的合法性基础和特定的制度逻辑。这里使用的“制度化”概念来自组织社会学中的新制度学派（Meyer \& Rowan，1977；参见周雪光，2003）。“制度化”的意义是，某个特定的组织行为或组织结构被社会或组织环境中有关各方所广泛接受，因此具有合法性或合理性；而组织环境中的制度化机制促就这类现象的稳定存在、重复再生。从这个角度来看，基层政府间的共谋现象恰具这些制度化特点：一方面，这些共谋行为大多与正式政策法令的明文规定不一致甚至相悖，旨在应付、欺骗、蒙蔽上级政府，因此多以非正式形式出现。但另一方面，这些行为并不是个别官员或个别部门的私人活动；在许多情形下，它们在正式组织权力结构下公开运作，以政府部门的组织权威辅以实施，甚至是大张旗鼓地加以部署安排。而这些做法已经成为上下级政府甚至中央政府的“共有常识”（common knowledge）。正是在这个意义上，我们说基层政府间共谋现象是特定的组织环境中制度化了的非正式行为。
\par\vspace{0.3em}{\footnotesize\color{gray}（周雪光. (2009). 基层政府间的“共谋现象”——一个政府行为的制度逻辑. 开放时代, (12), 40-55.）}
\end{document}